% Separate post-CP23 candidate; exact proof and certificate are audited separately.
\subsection{The canonical high quotient at \(p=5,e=12\)}

Throughout this subsection, after the explicitly indicated unramified
base change, $K$ is complete discretely valued of mixed characteristic
$(0,5)$ with algebraically closed residue field, $e_K=12$ and
$\mu_5\subset K$. Put $N=24$. For two wholly potentially supersingular
height-two groups, enlarge the common tame good field to $L/K$ of degree
$N$, and normalize $v_L$ to have value group $\mathbf Z$.
Then $v_L(5)=288$. In the nonsquare finite-determinant case a common
finite height-one twist makes that determinant primitive of order four.
Choose $\sigma\pi_L=\zeta\pi_L$ with $\zeta$ primitive of order $24$.
As in B53 and B60, the actual special tangent exponent $c$ satisfies
$c\equiv1\pmod4$, and its special character is the Frobenius orbit
$\{c,5c\}$ modulo $24$.

\begin{lemma}[Newton and point-field rules at index twelve]
\label{lem:B-RD65-Newton}
\status{Candidate; complete Newton and lower-break proof written.}
In the nonsquare normalized case, a reducible residual module has a
unique high-valuation line $H$ and an integer $B$ satisfying
\[
 0<B<60,\qquad B\equiv c\pmod6,\qquad
 \alpha=v_L(T(Q))=72-B\ (0\ne Q\in H),\quad
 \beta=v_L(T(Q))=B/5\ (Q\notin H).                    \tag{B-RD65.1}
\]
The high line is rational over $L$ and its tame character is
$\zeta^{c-B}$. Every point $P$ of order $25$ with $5P\notin H$
has valuation $B/125$.
For two such groups with isomorphic actual $25$-torsion:
\begin{enumerate}
\item If their residual wild image is nontrivial, their high lines
coincide. If neither $B_i$ is divisible by $5$, then $B_1=B_2$ and
$c_1=c_2$ modulo $24$. It is impossible that exactly one $B_i$ is
divisible by $5$.
\item If their residual wild image is trivial, write $B_i=5b_i$.
If both $b_i$ are prime to $5$, then $b_1=b_2$, their high lines
coincide, and $c_1=c_2$ modulo $24$.
\end{enumerate}
\end{lemma}
\begin{proof}
We repeat the calculations of B60--B61 with the actual new absolute
index, rather than applying their minimal-ramification statements
outside their scope. In the universal height-two coordinate,
$[5](T)$ has the three Newton vertices
$(1,288),(5,b),(25,0)$, where $b=v_L(u)$; the intermediate coefficient
bounds are exactly those established from Hopkins--Gross in
Lemma~\ref{lem:B-RD49-point}. If $b\ge240$, all nonzero $5$-torsion
points have valuation $12$. Their leading-coordinate map is injective,
so wild inertia acts trivially. The tame representation on its
leading plane is irreducible over $\mathbf F_5$ or scalar; the latter
would have square determinant, contrary to the normalization. Thus
reducibility implies $0<b<240$.

The high line has valuation $(288-b)/4$, and its character has values
in $\mathbf F_5^\times$, so it is killed by the tame good field $L$.
Its coordinate valuation is integral. Write $b=4B$; this gives
$\alpha=72-B$, $\beta=B/5$ and $0<B<60$. Tame averaging of the
formal parameter preserves the degree-five coefficient valuation
$b<288$. Compatibility with the exact action $T\mapsto\zeta^c\sigma(T)$
gives $b\equiv4c\pmod{24}$, hence $B\equiv c\pmod6$.
The high point's leading coefficient gives the character
$\zeta^{c+72-B}=\zeta^{c-B}$. In particular $B$ is odd.
For a target of valuation $\beta=B/5<12$, the inverse equation
$[5](T)=T(5P)$ has a single Newton segment of degree $25$, because
both $288$ and $4B$ exceed $\beta$. Thus $v_L(T(P))=B/125$.

Let $R=L(M[5])$. A reducible residual image over $L$ is either
trivial or cyclic unipotent of order five. In the latter case $H$
is its unique fixed line. For an element $g\ne1$ of
$\operatorname{Gal}(R/L)$ and a low point $Q$, one has
$v_R(T(Q))=B$ and
$v_R(T(gQ)-T(Q))=5\alpha=5(72-B)$.
The exact lower-break lemma~\ref{lem:B-RD60-break} therefore gives
\[
 r=360-6B\quad(5\nmid B),\qquad
 r<360-6B\quad(5\mid B),                              \tag{B-RD65.2}
\]
where $r$ is the single actual lower break of $R/L$.
Two prime-to-five $B_i$ are equal from the common break, and then
from the common high-line character their $c_i$ are equal.
If precisely one $B_i$ is divisible by five, index so that
$5\nmid B_1$ and $5\mid B_2$. Equation (B-RD65.2) forces $B_1>B_2$.

Choose corresponding $P_i$ with $5P_i$ outside the common high line.
Over $R$, their valuations are $B_i/25$. On the first side the
valuation denominator forces the point field $E=R(P_i)$ to have
degree at least $25$, and the elementary-abelian point-field
Lemma~\ref{lem:B-RD61-point-field} bounds it by $25$.
Thus the common field is Galois of degree $25$, every $5$-torsion
translation occurs, $[E:L]=125$, and $v_E(T_i(P_i))=B_i$.
Choose a translation outside $H$. Its difference has valuation
$25B_i$ in model $i$. The first coordinate has prime-to-five
valuation, so its actual lower break is $24B_1$. Applying only the
lower-bound part of the same lemma to the second coordinate gives
$25B_2\ge B_2+24B_1$, contradicting $B_1>B_2$.

In the tame residual case $R=L$ and $B_i=5b_i$ is necessary for
the low point's valuation to be integral. Choose a common residual
vector outside both high lines and lift it to $P_i$ of order $25$.
Their valuations are $b_i/25$. If both $b_i$ are prime to five,
the common point field has degree $25$ over $L$ and both normalized
coordinate valuations equal $b_i$. A translation outside both high
lines gives the same actual break $24b_1=24b_2$. Thus $b_1=b_2=b$.
If their high lines differed, a translation high only in the first
model would have respective difference valuations $25\alpha_1$
and $25b$, while the two coordinates have the same prime-to-five
valuation $b$. The equality part of the break lemma would force
$\alpha_1=b$, contradicting the strict high/low inequality.
Their high lines therefore agree, and their characters give $c_1=c_2$.
\end{proof}

\begin{lemma}[The exact canonical quotient transformation]
\label{lem:B-RD65-quotient}
\status{Candidate; complete integral quotient and Newton proof written.}
Let a group in Lemma~\ref{lem:B-RD65-Newton} have parameter $B\ne12$,
and let $M'=M/H$ be its quotient by the high line. It is again wholly
potentially supersingular. Its unequal Newton parameter is
\[
 B'=\begin{cases}5B,&B<12,\\72-B,&B>12.\end{cases}      \tag{B-RD65.3}
\]
Its special character is the same as that of $M$.
If two actual $125$-torsion modules are isomorphic and the isomorphism
preserves their high lines, then the actual $25$-torsion modules of
their high quotients are isomorphic.
\end{lemma}
\begin{proof}
Over the good DVR take the finite flat closure of $H$. Its quotient
is a good height-two formal group, isogenous to the original group.
Choose integral parameters for the source and quotient. The quotient
map has special degree five and, by Weierstrass preparation, its
parameter series is a unit times
$\prod_{Q\in H}(T-T(Q))$. The unit has valuation zero at every
positive-valuation point. In particular, if $v_L(T(P))<\alpha$,
then the quotient coordinate of $P$ has valuation $5v_L(T(P))$.
This uses the actual integral quotient, without a scaling of its
uniformizer or an identification of its Hasse coefficient with
that of the source.

The image of $M[5]/H$ is a line of $M'[5]$ with nonzero coordinate
valuation $5\beta=B$. All other points of $M'[5]$ come from points
$P$ with $5P$ a nonzero member of $H$.
If $B>12$, the inverse Newton equation with target valuation
$\alpha=72-B$ has one segment: the condition is
$4B>4\alpha/5$, equivalent to $B>12$.
Every preimage has valuation $\alpha/25<\alpha$, so the other
$20$ quotient points have valuation $\alpha/5$.
Since $B>\alpha/5$, the high and low valuations in the quotient
are $B$ and $(72-B)/5$, which gives $B'=72-B$.

If $B<12$, that inverse equation has five roots of valuation
$(\alpha-4B)/5=(72-5B)/5$ and twenty roots of valuation $B/5$.
Both valuations are less than $\alpha$. After quotienting by $H$,
the high preimages give four points of valuation $72-5B$, and
all other nonzero quotient points have valuation $B$.
(The latter include the four points of $M[5]/H$.) Hence $B'=5B$.
These counts follow by taking all four nonzero target points of $H$
and dividing the preimage counts by five. They also show that the
claimed unequal valuations are strict. The excluded value $B=12$
cannot in fact occur under $B\equiv c\pmod6$, $c\equiv1\pmod4$.

The special isogeny identifies the rational Dieudonne modules
with their descent actions, so the requested special characters
are unchanged. More concretely its special kernel is the unique
rank-five subgroup, and its target is the Frobenius twist of the
special group; the new tangent exponent is $5c$ modulo $24$, which
has the same unordered Frobenius orbit.
Finally the inverse image of $M'[25]$ is the actual finite subgroup
\[
       \{P\in M[125]:25P\in H\}.
\]
Dividing it by $H$ gives $M'[25]$ as a finite generic group scheme.
A $125$-torsion isomorphism preserving $H$ therefore induces the
asserted $25$-torsion isomorphism. This is the exact quotient
construction, not a comparison of isogeny classes of lattices.
\end{proof}

\begin{proposition}[All nonsquare formal comparisons at index twelve]
\label{prop:B-RD65-nonsquare}
\status{Candidate; complete residual-case proof written.}
Two normalized potentially supersingular height-two groups over $K$
with isomorphic actual $125$-torsion have the same finite special
inertia character.
\end{proposition}
\begin{proof}
First consider an irreducible residual module. An unequal Newton
polygon would provide an invariant high line, so both polygons
have one slope. The leading-coordinate calculation gives residual
tame eigenvalue orbit $\{c+12,5(c+12)\}$ modulo $24$.
As $5\cdot12\equiv12\pmod{24}$, equality of these orbits is exactly
equality of the special orbits $\{c,5c\}$.
A one-slope reducible representation would be scalar and have
square determinant, already excluded. It remains to consider the
unequal polygons of Lemma~\ref{lem:B-RD65-Newton}.

If residual wild inertia is nontrivial, its fixed high line is
common. That lemma handles the cases where either $B_i$ is prime
to five. If both are divisible, then $B_i=5b_i$ with
\[
                       b_i\in\{1,3,5,7,9,11\}.        \tag{B-RD65.4}
\]
Quotient by the common high line. Their $25$-torsion modules remain
isomorphic by Lemma~\ref{lem:B-RD65-quotient}.
For every $B_i\ge15$, the new parameter $72-B_i$ is prime to five.
The only smaller possibility is $B_i=5$, for which $B_i'=25$.
If both original parameters are $5$, their common high-line
character already identifies $c_i$. If both are at least $15$,
both quotient parameters are prime to five; their nonintegral low
valuations force nontrivial residual wild inertia in the quotients.
Lemma~\ref{lem:B-RD65-Newton} makes their parameters and special
characters equal. If exactly one original parameter is $5$, the
quotients have one divisible and one prime-to-five parameter, so
the same lemma forbids their $25$-torsion isomorphism.
Isogeny invariance transfers the conclusion to the original pair.

Suppose instead residual inertia is tame. Again write $B_i=5b_i$
with the list (B-RD65.4). If neither $b_i$ is divisible by five,
Lemma~\ref{lem:B-RD65-Newton} identifies the characters.
The only exceptional value is $b=5$, namely $B=25$.
If both parameters have this value, their allowed tangent exponents
are $1$ or $13$ modulo $24$. Their high-line characters are respectively
$1$ or $-1$, and their common finite determinant is primitive of
order four, say $\eta$. The residual unordered characters are
therefore $\{1,\eta\}$ and $\{-1,-\eta\}$, respectively. These sets
are distinct, so a common residual module identifies the special
character.

Finally suppose $b_2=5$ and $b_1=b$ is prime to five. Choose a common
residual vector outside both high lines and lift it to $P_i$ of
order $25$. The common point field $E=L(P_i)$ has degree $25$,
forced by the first valuation $b/25$, and all translations occur.
The two coordinates have normalized valuations $b$ and $5$.
If the high lines differed, choose a translation high in model one
and low in model two. The first coordinate's prime-to-five valuation
gives lower break $25(72-5b)-b$, whereas the lower-bound inequality
for the second coordinate would require
\[
                 125\ge5+25(72-5b)-b.                 \tag{B-RD65.5}
\]
This is impossible for every $b\in\{1,3,7,9,11\}$: its right side
is at least $419$.
Thus the high lines agree. A translation outside that line has
break $24b$ in the first coordinate, so the second coordinate gives
$125\ge5+24b$. Hence $b\le5$, leaving $b=1$ or $3$.
Their high-line characters now give
\[
                        c_1-c_2\equiv5b-25\pmod{24}.
                                                        \tag{B-RD65.6}
\]
Since both $c_i\equiv1\pmod4$, this forces $b\equiv1\pmod4$,
so only $b=1$ remains. Then $c_2=1$ or $13$, while
$c_1=c_2+4$ is $5$ or $17$, respectively. These are the same
Frobenius orbits. This exhausts the tame case.
\end{proof}

\begin{lemma}[The square-determinant formal comparison at index twelve]
\label{lem:B-RD65-square}
\status{Candidate; complete actual-target two-Frobenius proof written.}
Two wholly potentially supersingular height-two groups over $K$
with common square finite determinant and isomorphic actual
$125$-torsion have the same finite special inertia character.
\end{lemma}
\begin{proof}
Twist both groups by the same finite square root to make their
full determinants cyclotomic. Their finite special eigenvalues
then have norm one, so their common faithful good field is tame
of degree $m\mid6$, as in Proposition~\ref{prop:B-RD50-square}.
Its absolute index is $a=12m\le72$.
If the special characters differ, their two Frobenius eigenvalue
orbits are disjoint. The integral tame projector and the actual
perfect pairings of Proposition~\ref{prop:B-RD64-formal} make the
special graph closure in the $125$-torsion equal to
$\ker F_M^3\times\ker F_N^3$. The actual quotient $Q$ has special
fibre $M_0^{(5^3)}\times N_0^{(5^3)}$.
The map $M\times N\to Q\times M$ used there has no nonzero finite
flat kernel and has special kernel $0\times\ker F_N^3$.

We check the required logarithm bound on its actual target
$G=Q\times M$, without any transferred source thickness.
The invariant-differential identity makes all nonlinear coefficients
of total degree below five divisible by five; the product-height-two
special fibre makes coefficients of degrees five through twenty-four
lie in the maximal ideal. Starting with an augmentation vector of
valuation at least one, two applications of $[5]$ therefore give
coefficient valuations at least
\[
 q_1\ge6,\qquad q_2\ge\min\{a+6,31\}\ge a/4,
                    \qquad 12\le a\le72.                  \tag{B-RD65.7}
\]
The degree-at-least-twenty-five terms are larger throughout.
The closed coefficientwise logarithm estimate and the exact
Cartier bounded-projection/saturation argument of
Lemma~\ref{lem:B-RD64-weighted-kernel} now apply with two
Cartier Frobenius operators. Only its coefficient estimate has
been replaced by (B-RD65.7); its integral module argument is
unchanged. They force $F^2\ker f_k=0$, whereas
$F^2\ker F_N^3$ has rank five. This contradiction proves the lemma.
\end{proof}

\begin{theorem}[The next five-adic least index]
\label{thm:B-RD65}
\status{Candidate; complete full-monodromy reduction written.}
In the full B.4 monodromy setting, for height two, complete mixed
characteristic $(0,5)$ fields with perfect residue, $\mu_5\subset K$
and $e_K=12$, the exact least manuscript index is
\[
                           \nu_{\min}(5,12,2)=2.       \tag{B-RD65.8}
\]
\end{theorem}
\begin{proof}
The counterexample of Theorem~\ref{thm:B-RD63-arc} at $n=2$
exists because $12>5$, so index one need not suffice.
For an actual $125$-torsion comparison, pass to the completed
maximal unramified base. As $5\nmid12$, the cyclotomic image
modulo $125$ has order $25$. The determinant argument of
Theorem~\ref{thm:B-RD64} identifies the total dimension and the
finite determinant. Dimensions zero and two, and comparisons of
two filtered height-one systems, are determined by generic
mod-five inertia as proved there. A mixed whole-supersingular
versus filtered comparison is impossible by
Proposition~\ref{prop:B-RD63-mixed}, since $12\le5^3$.
For two whole potentially supersingular groups, the square
finite-determinant case is Lemma~\ref{lem:B-RD65-square}.
In the nonsquare case a common finite twist gives exactly the
normalization used above, and Proposition~\ref{prop:B-RD65-nonsquare}
identifies the characters. Every allowed comparison is covered,
so index two suffices and is least.
\end{proof}

\begin{remark}[Scope and retained finite checks]
\label{rem:B-RD65-scope}
\status{Candidate scope; no other least-index assertion.}
This argument uses the actual high-line quotient and the exact
$125$-to-$25$ subgroup quotient. It does not infer a rank-one lift
from a small Hasse valuation, and does not transfer thickness to
an isogenous quotient. The Newton and point-field rules are proved
at the new absolute index $288$. The finite supplement checks
only their displayed arithmetic and case exhaustion; it is not
used as a substitute for the formal-group or Galois proofs.
No exact least index at $p=5,e=16$ or at other new ramification
indices is asserted here.
The exact supplement
\path{code/reduction_descent/verify_p5_e12_high_quotient.py}
checks $60$ normalized Newton records, $120$ inverse polygons,
$60$ quotient point counts, $4716$ common-break rows (including
$4524$ divisible-parameter rows), $72$ tame residual pairs,
$25$ exceptional inequalities, $36$ one-slope orbit comparisons
and four square-determinant logarithm bounds. The author run passed
in $0.008917200997530017$ seconds at 20:33:10 UTC.
\end{remark}
